% 乘积可测结构：(a) A=[0.4,1.4]∪[2.0,2.8]、B=[0.5,1.5] 的可测矩形 A×B 由两块组成；(b) 单位圆盘四分之一（半径 2）由 0.5 网格的蓝色矩形与 0.25 网格补充的绿色矩形内逼近（各矩形取在圆内的最大网格高度），矩形可数并生成非矩形集合。
\begin{tikzpicture}[>=Stealth,
  ax/.style={draw=FigureStroke,line width=.6pt,->},
  rect/.style={draw=FigureBlue,fill=FigureBlueFill,line width=.7pt},
  proj/.style={draw=FigureGreen,line width=2pt},
  note/.style={font=\FigureNoteFont\fontsize{7.5}{9}\selectfont,text=FigureMuted}]
\begin{scope}
  \draw[ax] (-.1,0)--(3.3,0) node[right] {$X$};
  \draw[ax] (0,-.1)--(0,2.1) node[above] {$Y$};
  \draw[rect] (.4,.5) rectangle (1.4,1.5);
  \draw[rect] (2.0,.5) rectangle (2.8,1.5);
  \draw[proj] (.4,0)--(1.4,0); \draw[proj] (2.0,0)--(2.8,0);
  \draw[proj] (0,.5)--(0,1.5);
  \node[below] at (1.6,-.08) {$A$};
  \node[left] at (-.05,1.0) {$B$};
  \node at (.9,1.0) {$A\times B$};
  \draw[FigureMuted,dashed,line width=.4pt] (.4,0)--(.4,.5) (1.4,0)--(1.4,.5) (2.0,0)--(2.0,.5) (2.8,0)--(2.8,.5);
  \draw[FigureMuted,dashed,line width=.4pt] (0,.5)--(.4,.5) (0,1.5)--(.4,1.5);
  \node[align=center] at (1.6,-.85) {(a) 可测矩形};
\end{scope}
\begin{scope}[xshift=4.6cm]
  \draw[ax] (-.1,0)--(2.5,0) node[right] {$x$};
  \draw[ax] (0,-.1)--(0,2.4) node[above] {$y$};
  \foreach \i/\h in {0/1.5,1/1.5,2/1.0}{
    \draw[rect] ({\i*.5},0) rectangle ({\i*.5+.5},\h);}
  \foreach \a/\lo/\hi in {0/1.5/1.75,.25/1.5/1.75,.5/1.5/1.75,1/1.0/1.5,1.25/1.0/1.25,1.5/0/.75}{
    \draw[draw=FigureGreen,fill=FigureGreenFill,line width=.6pt] (\a,\lo) rectangle (\a+.25,\hi);}
  \draw[FigureStroke,line width=.9pt,domain=0:90,samples=60] plot ({2*cos(\x)},{2*sin(\x)});
  \node[note,anchor=west] at (1.55,1.95) {四分之一圆盘};
  \node[align=center] at (1.1,-.85) {(b) 矩形的不交并逼近};
\end{scope}
\end{tikzpicture}
